\begin{abstract}
Speech recognisers are scored by one overall error rate. Yet in English--Mandarin code-switched speech, the few words from the other language (\emph{embedded} words) often carry the meaning, and they may be the first to suffer when a model is shrunk to run on a device. We split Whisper's errors by each word's role in its sentence, on two corpora (ASCEND and NSC Part~4) and 120 recordings of our own, and find three things. First, the same word is wrong more often when it is embedded; fine-tuning removes this penalty for Mandarin but not for English. Second, two training losses aimed at embedded words, and language-ID re-ranking applied only to clips with a switch, give small but statistically reliable gains (best system: 20.6\% error on NSC, down from 55.1\%). Third, 8-bit quantization causes no detectable harm, while 4-bit hurts embedded words more than other words in the smaller models; our repairs recover about half to four-fifths of that loss in NF4, but not in the format used on device. We therefore recommend 8-bit models for on-device use.
\end{abstract}
